\kchapter{StatefulSets, DaemonSets, Jobs, and CronJobs}{sts}

\begin{objectives}
\begin{itemize}
  \item which workload controller to choose for which job;
  \item StatefulSets for databases and clustered software;
  \item DaemonSets for one-Pod-per-node agents;
  \item Jobs and CronJobs for batch and scheduled work.
\end{itemize}
\end{objectives}

\section{Choosing a workload type}

\begin{center}
\ktablestart
\begin{tabularx}{\linewidth}{P{3.3cm} L L}
\kheader \kth{Controller} & \kth{Use when\ldots} & \kth{Examples} \\
\kicon{deploy}\ Deployment & Pods are interchangeable and stateless. & Web
   apps, APIs, workers reading a queue. \\
\kicon{sts}\ StatefulSet & Each Pod needs a stable name and its own disk. &
   PostgreSQL, Kafka, Elasticsearch, ZooKeeper, Redis cluster. \\
\kicon{ds}\ DaemonSet & You need exactly one Pod on every (or some) node. &
   Log collectors, node monitoring agents, CNI and storage plugins. \\
\kicon{job}\ Job & Work runs to completion, then stops. & DB migrations,
   batch processing, report generation. \\
\kicon{cronjob}\ CronJob & A Job on a schedule. & Nightly backups, cleanup,
   periodic sync. \\
\end{tabularx}
\end{center}

\section{StatefulSets}

Deployments treat Pods as cattle: random names, shared storage (if any), any
order. Stateful systems need more. A \kterm{StatefulSet} guarantees:

\begin{itemize}
  \item \textbf{stable, ordinal names}: \kc{db-0}, \kc{db-1}, \kc{db-2} --- a
        replaced \kc{db-1} is still \kc{db-1};
  \item \textbf{stable network identity} through a headless Service:
        \kc{db-0.db.default.svc.cluster.local};
  \item \textbf{stable storage}: each Pod gets its own PersistentVolumeClaim
        from \kc{volumeClaimTemplates}, which follows it across rescheduling;
  \item \textbf{ordered} start-up, scaling, and rolling updates (0, 1, 2 \ldots).
\end{itemize}

\begin{center}
\begin{tikzpicture}[
  p/.style={fill=k8sblue, text=white, rounded corners=4pt, minimum width=2.1cm,
            minimum height=0.9cm, font=\small\bfseries},
  d/.style={cylinder, shape border rotate=90, aspect=0.25, draw=k8sdeep,
            fill=k8slight, minimum width=1.4cm, minimum height=1.0cm, font=\scriptsize},
  >={Stealth[length=2.5mm]}]
  \foreach \i in {0,1,2}{
    \node[p] (p\i) at (\i*3.2,0) {\kicon[0.9em]{pod}\ db-\i};
    \node[d] (d\i) at (\i*3.2,-1.6) {data-db-\i};
    \draw[<->, thick, k8sdeep] (p\i) -- (d\i);
  }
  \draw[->, thick, k8sorange] (p0.east) -- (p1.west);
  \draw[->, thick, k8sorange] (p1.east) -- (p2.west);
  \node[font=\scriptsize, text=k8sorange] at (4.8,0.75) {created in order};
\end{tikzpicture}
\end{center}

\begin{yamlcode}{postgres-statefulset.yaml}
apiVersion: v1
kind: Service
metadata:
  name: db
spec:
  clusterIP: None              # "headless": DNS returns Pod IPs directly
  selector:
    app: db
  ports:
    - port: 5432
---
apiVersion: apps/v1
kind: StatefulSet
metadata:
  name: db
spec:
  serviceName: db              # the headless Service above
  replicas: 1
  selector:
    matchLabels:
      app: db
  template:
    metadata:
      labels:
        app: db
    spec:
      containers:
        - name: postgres
          image: postgres:16-alpine
          env:
            - name: POSTGRES_PASSWORD
              valueFrom:
                secretKeyRef: { name: db-secret, key: password }
            - name: PGDATA
              value: /var/lib/postgresql/data/pgdata
          ports:
            - containerPort: 5432
          volumeMounts:
            - name: data
              mountPath: /var/lib/postgresql/data
  volumeClaimTemplates:        # one PVC per Pod: data-db-0, data-db-1, ...
    - metadata:
        name: data
      spec:
        accessModes: ["ReadWriteOnce"]
        resources:
          requests:
            storage: 5Gi
\end{yamlcode}

\begin{warningbox}
A StatefulSet gives you stable building blocks; it does \textbf{not} make a
database highly available, replicate data, or take backups. For production
databases use a managed service or a mature \textbf{Operator} (e.g.\
CloudNativePG, Percona, Strimzi for Kafka) --- see Chapter 19.
\end{warningbox}

Deleting a StatefulSet does \emph{not} delete its PVCs by default --- your data
survives. Control this with \kc{persistentVolumeClaimRetentionPolicy}.

\section{DaemonSets}

A \kterm{DaemonSet} runs one copy of a Pod on every node. When a node joins,
it gets the Pod; when a node leaves, the Pod is garbage-collected.

\begin{yamlcode}{node-agent-daemonset.yaml}
apiVersion: apps/v1
kind: DaemonSet
metadata:
  name: node-agent
  namespace: monitoring
spec:
  selector:
    matchLabels:
      app: node-agent
  template:
    metadata:
      labels:
        app: node-agent
    spec:
      tolerations:                         # also run on control-plane nodes
        - key: node-role.kubernetes.io/control-plane
          operator: Exists
          effect: NoSchedule
      containers:
        - name: agent
          image: prom/node-exporter:v1.8.2
          ports:
            - containerPort: 9100
          resources:
            requests: { cpu: 20m, memory: 32Mi }
            limits:   { memory: 64Mi }
\end{yamlcode}

Use a \kc{nodeSelector} or node affinity in the template to target only some
nodes (for example, only GPU nodes).

\section{Jobs}

A \kterm{Job} creates Pods that run until they \emph{succeed}. If a Pod fails,
the Job retries it (up to \kc{backoffLimit}).

\begin{yamlcode}{pi-job.yaml}
apiVersion: batch/v1
kind: Job
metadata:
  name: pi
spec:
  completions: 5             # need 5 successful Pods in total
  parallelism: 2             # run up to 2 at a time
  backoffLimit: 3            # retries before the Job is marked Failed
  activeDeadlineSeconds: 600 # give up after 10 minutes
  ttlSecondsAfterFinished: 3600   # auto-delete 1 hour after finishing
  template:
    spec:
      restartPolicy: Never   # Jobs must use Never or OnFailure
      containers:
        - name: pi
          image: perl:5.40
          command: ["perl", "-Mbignum=bpi", "-wle", "print bpi(500)"]
\end{yamlcode}

\begin{bashcode}{terminal}
kubectl apply -f pi-job.yaml
kubectl get jobs -w
kubectl logs job/pi
\end{bashcode}

\begin{protip}
\kc{completionMode: Indexed} gives each Pod an index (\kc{JOB\_COMPLETION\_INDEX}
env var), which is perfect for splitting work into shards: Pod 0 processes
file 0, Pod 1 file 1, and so on.
\end{protip}

\section{CronJobs}

A \kterm{CronJob} creates a Job on a schedule written in classic cron syntax.

\begin{center}
\begin{tikzpicture}[
  f/.style={fill=k8sblue, text=white, rounded corners=4pt, minimum width=1.5cm,
            minimum height=1cm, font=\ttfamily\Large\bfseries},
  l/.style={font=\scriptsize, text=k8snavy, align=center, text width=1.9cm}]
  \foreach \x/\v/\lbl in {0/30/minute\\(0-59), 2.2/2/hour\\(0-23), 4.4/*/day of\\month (1-31),
                          6.6/*/month\\(1-12), 8.8/1-5/{day of week\\(0-6; Sun=0)}}{
    \node[f] at (\x,0) {\v};
    \node[l] at (\x,-1.05) {\lbl};
  }
  \node[font=\footnotesize\itshape, text=k8sgray] at (4.4,-1.9) {``30 2 * * 1-5'' = 02:30 every Monday to Friday};
\end{tikzpicture}
\end{center}

\begin{yamlcode}{backup-cronjob.yaml}
apiVersion: batch/v1
kind: CronJob
metadata:
  name: nightly-backup
spec:
  schedule: "30 2 * * *"           # 02:30 every day
  timeZone: "Asia/Dhaka"           # optional; defaults to the controller's zone (UTC)
  concurrencyPolicy: Forbid         # skip if the previous run is still going
  successfulJobsHistoryLimit: 3
  failedJobsHistoryLimit: 3
  startingDeadlineSeconds: 300
  jobTemplate:
    spec:
      backoffLimit: 2
      template:
        spec:
          restartPolicy: OnFailure
          containers:
            - name: backup
              image: postgres:16-alpine
              command: ["sh", "-c", "pg_dump -h db -U postgres app > /backup/app-$(date +%F).sql"]
              env:
                - name: PGPASSWORD
                  valueFrom:
                    secretKeyRef: { name: db-secret, key: password }
              volumeMounts:
                - name: backup
                  mountPath: /backup
          volumes:
            - name: backup
              persistentVolumeClaim:
                claimName: backups
\end{yamlcode}

\begin{bashcode}{terminal}
kubectl get cronjobs
kubectl create job --from=cronjob/nightly-backup manual-run-1   # run it now
\end{bashcode}

\begin{recap}
\begin{itemize}
  \item Deployment: stateless and interchangeable. StatefulSet: identity + own disk.
  \item DaemonSet: one per node. Job: run to completion. CronJob: Jobs on a schedule.
  \item Stateful workloads in production usually deserve an Operator.
\end{itemize}
\end{recap}
